\documentclass[11pt,a4paper]{article}
\usepackage[margin=2.2cm]{geometry}
\usepackage{amsmath,amssymb,xcolor,enumitem,booktabs,graphicx}
\usepackage[normalem]{ulem}
\setlength{\parskip}{0.6em}\setlength{\parindent}{0pt}
\definecolor{newc}{HTML}{1f6f43}
\newcommand{\new}[1]{\textcolor{newc}{#1}}
\newcommand{\old}[1]{\textcolor{gray}{\sout{#1}}}
\newcommand{\note}[1]{{\small\itshape [#1]}}
\newcommand{\lab}[1]{\textsf{\footnotesize #1}}
\renewcommand{\ref}[1]{\lab{#1}}
\renewcommand{\eqref}[1]{(\lab{#1})}
\newcommand{\citep}[1]{[\lab{#1}]}
\newcommand{\citet}[1]{\lab{#1}}
\title{Proposed change, ``Choosing together''}
\date{24 November 2026, version 1}
\begin{document}\sloppy
\maketitle

\section*{How to read this}
This answers \texttt{method\_controllers\_comments.md} for
\lab{sec:method:controllers} in \texttt{4.Implementation.tex}. B2, B3, B4, C1
and E1 were already handled in [A119] and [A120] and stay as they are. This
change does D1, D2 and D4, and, as agreed on 24 November, B1 (one sentence
in the Discussion, the controller table kept) and D3 (the closing paragraph
dropped). Nothing in the model, the code or the figures changes.

Black text stays, \new{green} is new and \old{struck} text goes.
\note{Notes in brackets} say why. Nothing below is in the thesis yet.

\section*{Part 1. What the Problem Formulation already says}
\lab{sec:pf:controllers}, checked today:
\begin{itemize}[nosep]
  \item ``Centralised'': one network run per agent in priority order
    (``One decision runs the network only once per agent''), and the rule
    that an agent keeps the vertex and swap rules with the agents before
    it and may not move onto an agent that has not chosen yet.
  \item ``Section-based'': the zone takes over an agent that crosses into it,
    and ``with a single zone, the section-based controller is the
    centralised one''.
  \item ``What the controllers share'': they differ in how many agents one
    decision covers, and three things follow, what one decision sees,
    whether an agent knows the moves chosen before it, and whose reward an
    action is judged by.
\end{itemize}

\section*{Part 2. The section as it would read}

The centralised and section-based controllers use the same network, but
on a larger window and for one agent after another. The Problem Formulation
gives the order and the rule that keeps the agents of one decision apart
(Section~\ref{sec:pf:controllers}).

\textbf{Centralised.} The network runs on every vertex an agent can stand
on, once for each agent in priority order. Each run uses the agent's own
distance label and the vertices the earlier agents have chosen, and the move
it draws becomes part of what the next agent sees. \old{A timestep therefore
takes $m$ runs, one after another, so runtime grows with the size of the
fleet.} \new{The runs follow one another, so runtime grows with the size of
the fleet.} One pass shared by all agents, as in the single-agent setup of
\citet{vanKnippenberg2021}, would be cheaper, but then no agent would see
another's choice. \note{D4. One run per agent is in the Problem Formulation.}

\textbf{The claim mask.} The rule that keeps the agents of one decision
apart enters as a second mask\old{. It switches off a move to $w$ when an
earlier agent has chosen $w$, when $w$ holds an agent that has not chosen
yet, or when the earlier agent standing on $w$ has chosen to move to $v$,
which would be a swap.}\new{, the \emph{claim mask}
$M^{\prec}_i$~\eqref{eq:claimmask}. It is zero for a decentralised agent.}
A swap never arises inside one decision, because the earlier agent could not
move onto an agent that had not chosen yet. The swap rule does its work in
conflict resolution, where decentralised agents choose at the same time
(Figure~\ref{fig:claimmask}). Waiting is never switched off, so an agent
always has an action left. \old{This is the \emph{claim mask}
$M^{\prec}_i$~\eqref{eq:claimmask}, and it is zero for a decentralised
agent.} Conflict resolution therefore never overrides a centralised agent.
\note{D1. The three cases are the rule of the Problem Formulation. The mask
is named where it is introduced, and the vertex $v$ no longer appears
unnamed.}

\textbf{What the larger window adds.} \emph{[Unchanged, with its pooled
readout \lab{eq:pooledreadout}. It already says what a centralised agent
sees in detail and what only as an average (B1).]}

\textbf{Section-based.} Each zone does the same on its own part of the
warehouse. The zones are blocks of neighbouring aisles, as equal in size as
the aisles allow. A zone is not a category of SKUs, and the two are drawn
apart. Zone $Y_q$ runs the network for each of its agents in priority order,
on its own vertices together with a \emph{halo} $H_q$~\eqref{eq:halo}. The
halo is the strip of vertices outside the zone within field-of-view range of
it. The idea comes from distributed training of graph neural networks, where
each part of a split graph keeps copies of the neighbouring vertices it
needs, its halo vertices \citep{zheng2020distdgl}. Here the halo is as wide
as the field of view. After $L$ rounds of message passing a vertex has only
taken in the vertices within $L$ edges \citep{gilmer2017mpnn}, and the halo
is wider than the $L = 2$ rounds. So the network of a zone never reads past
its halo. Its vertices carry the same facts as any other vertex. No agent on
them chooses in this zone, and they are left out of the mean, because they
are only context. So an agent at the zone edge still sees its whole field of
view, and its features, crowding among them, mean the same as under the
other controllers.

\new{\textbf{At the border.}} A zone does not know the choices made in other
zones, so an agent treats every occupied halo vertex as holding an agent that
has not chosen yet, and does not move onto it. The weights are shared, so
moving into the next zone changes the agent's context but not its policy. The
zones run in parallel, and two zones may still send agents onto the same
empty vertex at the border. Conflict resolution removes these clashes as it
does everywhere else. \old{With $p=1$ the halo is empty and the
section-based controller is the centralised one.} So, unlike a centralised
agent, a section-based agent can be overridden at a zone border. This is
part of what deciding in zones costs, and the comparison keeps it.
\note{D2. The paragraph splits in two, the zone and its halo, then the
border. The $p=1$ sentence repeats the Problem Formulation. It was added
back in [A123] from the reviewer's green sentence, which missed that.}

[Figure \lab{fig:claimmask}, unchanged.]

\old{The three learned controllers share the network, the masks, the reward
of each agent and the training. They differ only in which agents choose
together, and so in the window the network runs on, in which earlier choices
an agent sees, and in the credit each agent gets. This is the contrast RQ1
asks about.} \note{D3. The Problem Formulation's ``What the controllers
share'' says the same, with the same three differences.}

\section*{Part 3. The Discussion}
\texttt{5.Experiments.tex}, ``What was expected'':

\dots\ The centralised controller did not gain the most. At 16 agents it
serves about as much as the others under congestion-aware storage. \new{A
possible reason is that its network sees the agents near it in detail and
the rest of the warehouse only as an average. It then differs from the
others mainly through the claim mask and the shared credit
(Section~\ref{sec:method:controllers}), not through seeing more.} Closing
the loop changed little for Token Passing, \dots

\note{B1. The controller table of the Problem Formulation keeps ``the whole
warehouse'', since it states what the controller may see. How much of it the
network uses is method.}

\section*{Questions}
None. Approved in chat on 24 November, apart from reading this page.

\end{document}
